%=====================================================================
\section{Solutions and hints}
\label{sec:solutions}

\paragraph{Exercise~\ref{ex:1-1}.}
$(A^\cl)^\dagger$: $(X,A^\cl Y)=\tfrac12\Tr[X^\dagger(AY+YA)]=\tfrac12\Tr[(AX+XA)^\dagger Y]=(A^\cl X,Y)$ using
$A=A^\dagger$ and cyclicity. $(A^\qu)^\dagger$: $(X,[A,Y])=\Tr[X^\dagger AY-X^\dagger YA]=\Tr[(AX-XA)^\dagger(-1)\cdots]$,
giving $-(A^\qu X,Y)$. Hence $iA^\qu$ is Hermitian and $\Lc=-iH^\qu$ is anti-Hermitian, as the
generator of a unitary (isometric) evolution in Liouville space must be.

\paragraph{Exercise~\ref{ex:1-2}.}
With $\mathrm{vec}$ in the basis $\{\ket0\bra0,\ket0\bra1,\ket1\bra0,\ket1\bra1\}$,
$\Lc_S=-i(\Id\otimes H_S-H_S^T\otimes\Id)=\operatorname{diag}(0,-i\omega_0,i\omega_0,0)$. The change of basis to
$\sigma_\mu/\sqrt2$ is the unitary with columns $\mathrm{vec}(\sigma_\mu)/\sqrt2$; it maps the
$2\times2$ block $\operatorname{diag}(-i\omega_0,i\omega_0)$ on $(\ket0\bra1,\ket1\bra0)$ to the real rotation generator
$\begin{psmallmatrix}0&-\omega_0\\ \omega_0&0\end{psmallmatrix}$ on $(\sx,\sy)$.

\paragraph{Exercise~\ref{ex:1-3}.}
Trace preservation fixes the first row; Hermiticity preservation makes the matrix
real. For a unital qubit map with $T=\operatorname{diag}(\lambda_1,\lambda_2,\lambda_3)$, positivity is
$|\lambda_i|\le1$ while complete positivity (Fujiwara--Algoet conditions) is
$|\lambda_1\pm\lambda_2|\le|1\pm\lambda_3|$; transposition $(1,-1,1)$ violates the latter.

\paragraph{Exercise~\ref{ex:1-4}.}
$(S^\cl)^2=\tfrac14(S^{+}+S^{-})^2=\tfrac14(2+2S^+S^-)$ using $S^2=\Id$: it maps $X\mapsto\tfrac12(X+SXS)$. Similarly
$(S^\qu)^2X=2X-2SXS$. Hence $(S^\cl)^2+\tfrac14(S^\qu)^2=1$ and $(S^\cl)^2-\tfrac14(S^\qu)^2=S^+S^-$, i.e.\
$X\mapsto SXS$. In matrix form, $(S^\cl)^2$ projects onto $\mathrm{span}(\Id,S)$ and $\tfrac14(S^\qu)^2$
onto its complement (the two Pauli directions orthogonal to $\hat n$).

\paragraph{Exercise~\ref{ex:1-5}.}
$[A^+,B^+]X=ABX-BAX=([A,B])^+X$; $[A^-,B^-]X=XBA-XAB=-([A,B])^-X$; $[A^+,B^-]X=AXB-AXB=0$.
For \cref{eq:product-clqu} insert $A^\pm=A^\cl\pm\tfrac12A^\qu$ into $(AB)^\cl=\tfrac12(A^+B^++B^-A^-)$:
$A^+B^+=A^\cl B^\cl+\tfrac12A^\cl B^\qu+\tfrac12A^\qu B^\cl+\tfrac14A^\qu B^\qu$ and
$B^-A^-=B^\cl A^\cl-\tfrac12B^\cl A^\qu-\tfrac12B^\qu A^\cl+\tfrac14B^\qu A^\qu$; adding and halving gives the
anticommutators and the two commutators, and subtracting gives $(AB)^\qu$. For
$A=\sx$, $B=\sy$ the commutator terms are essential: the two anticommutator terms cancel
each other, $\tfrac12\{\sx^\cl,\sy^\cl\}=-\tfrac18\{\sx^\qu,\sy^\qu\}$: the first swaps $\sx\leftrightarrow\sy$ with
weight $\tfrac12$ and annihilates $\Id$ and $\sz$, the second is its negative (check with the
matrices of Example~1.2), while $\tfrac14([\sx^\cl,\sy^\qu]+[\sx^\qu,\sy^\cl])=i\sz^\cl$ exactly, which is the
whole answer $(\sx\sy)^\cl=(i\sz)^\cl$. The script
\texttt{spin\_boson\_numerics.py} contains the vectorized superoperators; the identities
hold to machine precision for random $A,B$.

\paragraph{Exercise~\ref{ex:2-1}.}
Straightforward from $b_k(t)=b_k\ee^{-i\omega_kt}$, $\avg{b_k^\dagger b_k}=\bar n_k$, $\avg{b_kb_k^\dagger}=\bar n_k+1$.

\paragraph{Exercise~\ref{ex:2-2}.}
From \cref{eq:chi-def}, $\chi(\omega)=\tfrac2\pi\int_0^\infty\dd\omega'J(\omega')\,I(\omega,\omega')$ with
$I(\omega,\omega')=\int_0^\infty\dd t\,\ee^{i\omega t}\sin\omega't$. This integral does not converge in the
ordinary sense; it is defined as a distribution by a convergence factor $\ee^{-\epsilon t}$,
$\epsilon\to0^+$ (adiabatic switching; the one-sided range comes from the $\theta(t)$ in
$\chi$). \emph{Step 1.} For real $k$,
\begin{equation*}
  \int_0^\infty\!\dd t\,\ee^{ikt-\epsilon t}=\frac{1}{\epsilon-ik}=\frac{\epsilon}{k^2+\epsilon^2}+i\,\frac{k}{k^2+\epsilon^2}
  \ \xrightarrow{\epsilon\to0^+}\ \pi\delta(k)+i\,\mathrm P\frac1k ,
\end{equation*}
because the Lorentzian $\epsilon/(k^2+\epsilon^2)$ has width $\epsilon$ and area $\pi$, while
$k/(k^2+\epsilon^2)$ is odd and equals $1/k$ away from $k=0$ (Sokhotski--Plemelj). Compare the
two-sided integral $\int_{-\infty}^\infty\ee^{ikt}\dd t=2\pi\delta(k)$: the one-sided one is half the
delta plus the odd principal-value part that the other half would cancel.
\emph{Step 2.} Write $\sin\omega't=(\ee^{i\omega't}-\ee^{-i\omega't})/2i$ and apply Step 1 with
$k=\omega\pm\omega'$:
\begin{align*}
  I&=\frac{1}{2i}\Big[\pi\delta(\omega+\omega')-\pi\delta(\omega-\omega')+i\,\mathrm P\frac{1}{\omega+\omega'}-i\,\mathrm P\frac{1}{\omega-\omega'}\Big],\\
  \Re I&=\frac12\Big[\mathrm P\frac{1}{\omega+\omega'}-\mathrm P\frac{1}{\omega-\omega'}\Big]=-\,\mathrm P\frac{\omega'}{\omega^2-\omega'^2},\qquad
  \Im I=\frac\pi2\big[\delta(\omega-\omega')-\delta(\omega+\omega')\big].
\end{align*}
(Without complex exponentials: $\int_0^\infty\sin\omega't\,\sin\omega t\,\ee^{-\epsilon t}\dd t
=\tfrac12\big[\tfrac{\epsilon}{(\omega-\omega')^2+\epsilon^2}-\tfrac{\epsilon}{(\omega+\omega')^2+\epsilon^2}\big]$, two Lorentzians
that become the two deltas.) \emph{Step 3.} Insert into $\chi(\omega)$. The deltas are
integrated against $J(\omega')$ over $\omega'>0$, so for $\omega>0$ only $\delta(\omega-\omega')$ fires and
$\chi''(\omega)=J(\omega)$, while for $\omega<0$ only $\delta(\omega+\omega')$ fires, with its minus sign, and
$\chi''(\omega)=-J(-\omega)$: together $\chi''(\omega)=\sgn(\omega)J(|\omega|)$, odd in $\omega$ as a dissipative
response must be. The real part of the same formula gives
$\chi'(\omega)=\tfrac2\pi\,\mathrm P\!\int_0^\infty\dd\omega'\,\omega'J(\omega')/(\omega'^2-\omega^2)$, the Kramers--Kronig relation
of \cref{sec:lec2}: one regularized integral yields both the dissipative part (deltas) and
the reactive part (principal value), tied together by the causal $\theta(t)$. Two
remarks. The sign of $i\epsilon$ is the only place where retarded versus advanced is
decided: $\int_{-\infty}^0\ee^{ikt+\epsilon t}\dd t=\pi\delta(k)-i\mathrm P(1/k)$, same deltas, opposite
principal value, hence $\chi^A(\omega)=\chi^R(\omega)^*$. And the finite-time version,
$\Re\int_0^t\ee^{i(\omega-\omega_0)\tau}\dd\tau=\sin[(\omega-\omega_0)t]/(\omega-\omega_0)\to\pi\delta(\omega-\omega_0)$, is the filter of
\cref{eq:kofman} that turns the running rates of \cref{sec:lec6} into Markov rates: a sharp
cutoff in time instead of an exponential one.

For the finite-$T$ Ohmic
kernel use $\int_0^\infty\omega\coth(\beta\omega/2)\cos\omega t\,\dd\omega=-(\pi/\beta)^2/\sinh^2(\pi t/\beta)$ for $t\ne0$ (the
$t=0$ singularity carries the positive weight that makes $S(\omega)\ge0$). $S(t)<0$ for all
$t>0$ is compatible with $S(\omega)\ge0$ because $S(\omega)=\int S(t)\ee^{i\omega t}\dd t$ is dominated by the
short-time peak.

\paragraph{Exercise~\ref{ex:2-3}.}
The adjoint master equation gives $\dot a=-(i\Omega+\gamma/2)a$ for the Heisenberg operator
in the absence of the qubit; with the quantum regression theorem,
$\avg{a(t)a^\dagger}=\ee^{-(i\Omega+\gamma/2)t}$. Gaussianity: the Lindbladian is quadratic in
$a,a^\dagger$, so it maps Gaussian states to Gaussian states.

\paragraph{Exercise~\ref{ex:2-4}.}
$\Tr_B[(\cos\theta\sz+\sin\theta\sx)B\cdots(\cos\theta\sz+\sin\theta\sx)B]$ contains $\cos\theta\sin\theta\,C(t)$ cross
terms between $\sz$ and $\sx$ vertices. Two independent baths would give
$C_{zz}=\cos^2\theta C$, $C_{xx}=\sin^2\theta C$, $C_{zx}=0$. The cross terms are the $cs$ entries of
\cref{eq:K2-bloch}.

\paragraph{Exercise~\ref{ex:3-1}.}
Eight terms $\tilde\Lc_I^{a_1}\tilde\Lc_I^{a_2}\tilde\Lc_I^{a_3}$, $a_i=\pm$, each a contour-ordered
three-point function of $B$, which vanishes for a Gaussian state with $\avg B=0$.

\paragraph{Exercise~\ref{ex:3-2}.}
For $t>s$: $D^{\cl\qu}=\tfrac12[C(t-s)-C(s-t)+C(t-s)-C(s-t)]=C(t-s)-C(s-t)=\avg{[B(t),B(s)]}$. For $t<s$ the
later superoperator is $B^\qu(s)$ and the trace identity kills the term. Without the table:
write $D^{ab}$ as $\Tr_B[\text{later}\cdot\text{earlier}\cdot\rhoB]$; if the later one is $B^\qu$ the
trace vanishes; if both are $B^\cl$ the result is symmetric in $t,s$.

\paragraph{Exercise~\ref{ex:3-3}.}
$\Tr_S[\tilde S^\cl(t)\tilde S^\cl(s)\rho_S]=\tfrac12\avg{\{\tilde S(t),\tilde S(s)\}}=\cos^2\theta+\sin^2\theta\cos\omega_0(t-s)$ for the
diagonal state (independent of $r_z$), and
$\Tr_S[\tilde S^\cl(t)\tilde S^\qu(s)\rho_S]=\avg{[\tilde S(t),\tilde S(s)]}=-2i\sin^2\theta\,r_z\sin\omega_0(t-s)$ for $t>s$: the
qubit's response function, proportional to its polarization, as a two-level system's
susceptibility must be.

\paragraph{Exercise~\ref{ex:3-4}.}
See the text; $\chi''(\omega)=J(\omega)-J(-\omega)$ and $S(\omega)=J(\omega)+J(-\omega)$ give $S=\sgn(\omega)\chi''$ at $T=0$
only for $|\omega|$ where $J(-|\omega|)\ll J(|\omega|)$; the negative-frequency tail is the price of the
pseudomode.

\paragraph{Exercise~\ref{ex:4-1}.}
$\Tr_SK^{(2)}X=-\Tr_S[S^\qu(\cdots)]=0$ by the first trace identity for $S$. Hermiticity: with
$Y=\ee^{\Lc_S\tau}[S(\tau)S^\qu X-i\chi(\tau)S^\cl X]$, $Y^\dagger=\ee^{\Lc_S\tau}[-S(\tau)S^\qu X^\dagger+i\chi S^\cl X^\dagger]$, so
$(S^\qu Y)^\dagger=-S^\qu Y^\dagger=S^\qu\ee^{\Lc_S\tau}[S\,S^\qu-i\chi S^\cl]X^\dagger$ for Hermitian $X$.

\paragraph{Exercise~\ref{ex:4-2}.}
Input $\sz$: $S^\cl\sz=\cos\theta\,\Id$, $S^\qu\sz=-2i\sin\theta\,\sy$. Free rotation:
$\ee^{\Lc_S\tau}\sy=\cos\omega_0\tau\,\sy-\sin\omega_0\tau\,\sx$, $\ee^{\Lc_S\tau}\Id=\Id$. Then
$-S^\qu[\,S_\tau(-2i\sin\theta)(\cos\,\sy-\sin\,\sx)-i\chi_\tau\cos\theta\,\Id\,]$; use $S^\qu\Id=0$,
$S^\qu\sy=-2i\cos\theta\sx+2i\sin\theta\sz$, $S^\qu\sx=2i\cos\theta\sy$:
$=-S_\tau\sin\theta\,[-4\cos\theta\cos\,\sx+4\sin\theta\cos\,\sz-4\cos\theta\sin\,\sy]$, giving the last column of
\cref{eq:K2-bloch}. Input $\Id$: $S^\cl\Id=S$, $S^\qu\Id=0$; $\ee^{\Lc_S\tau}S=\cos\theta\sz+\sin\theta(\cos\,\sx+\sin\,\sy)$;
then $i\chi_\tau S^\qu[\cdots]=\chi_\tau[2cs(1-\cos)\sy+2cs\sin\,\sx-2s^2\sin\,\sz]$, the first column. The
$\sx,\sy$ columns are analogous. At $\theta=0$ only the $c^2$ block survives (dephasing of
$\sx,\sy$ at rate $4S_\tau$ with the rotation); at $\theta=\pi/2$ only the $s^2$ entries
(relaxation of $\sz$, dephasing of $\sy$ but not of $\sx$ in the lab frame, and the
$\chi$-driven inhomogeneity).

\paragraph{Exercise~\ref{ex:4-3}.}
The 15 pairings of six ordered points split into 10 irreducible and 5 reducible:
$(12)(34)(56)$, $(12)\times$\{2 irreducible 4-point\}, \{2 irreducible 4-point\}$\times(56)$.
(Irreducible chord diagrams: $1,1,2,10,74,\dots$)

\paragraph{Exercise~\ref{ex:4-4}.}
$\Gamma(t)=4\int_0^t\dd t_1\int_0^{t_1}\dd t_2\,g^2\ee^{-\gamma(t_1-t_2)/2}\cos\Omega(t_1-t_2)$. Short times:
$4g^2t^2/2=2g^2t^2$. Long times: $4t\int_0^\infty S(\tau)\dd\tau=2S(\omega{=}0)\,t$ with
$S(0)=2g^2(\gamma/2)/(\Omega^2+\gamma^2/4)$, the zero-frequency value of $J(\omega)+J(-\omega)$.

\paragraph{Exercise~\ref{ex:5-1}.}
$\Pc\Lc\Pc=\Lc_S\Pc$ uses $\Pc\Lc_I\Pc=0$; $\Pc\Lc\ee^{\Qc\Lc\tau}\Qc\Lc\Pc=\Pc\Lc_I\ee^{\Qc\Lc\tau}\Qc\Lc_I\Pc$ uses
$\Pc\Lc_0\Qc=0$. With $\avg B\ne0$, write $B=\avg B+\delta B$ and absorb $S\avg B$ into $H_S$.

\paragraph{Exercise~\ref{ex:5-2}.}
Secular reduction of \cref{eq:K2-bloch} with $\ee^{-\Lc_S\tau}$ acting on the input:
$\Gc_{xx}=-4c^2\int S$, $\Gc_{yy}=-4c^2\int S-4s^2\int S\cos\omega_0\tau$, $\Gc_{yx}=4s^2\int S\sin\omega_0\tau$,
$\Gc_{xy}=0$, $\Gc_{zz}=-4s^2\int S\cos\omega_0\tau$, $\Gc_{z\Id}=-2s^2\int\chi\sin\omega_0\tau$. Average $\Gc_{xx}$ and
$\Gc_{yy}$ (the difference oscillates at $2\omega_0$ in the rotating frame) to get $\Gamma_2$; the
precession frequency is $\sqrt{\omega_0(\omega_0+\Gc_{yx})}\simeq\omega_0+\tfrac12\Gc_{yx}$. The Lamb shift
contains $\int_0^\infty S\sin\omega_0\tau$ only: the $\chi$ terms sit in the first column, which is the
inhomogeneity, not a frequency. In the operator form, $\Delta(\omega)=\Im\int_0^\infty\ee^{i\omega s}C(s)\dd s=\int_0^\infty S\sin\omega s-\tfrac12\chi'(\omega)$
and $\delta\omega=\sin^2\theta[\Delta(\omega_0)-\Delta(-\omega_0)]$; $\chi'$ is even and cancels.

\paragraph{Exercise~\ref{ex:5-3}.}
Secular: $\Gc=-i[\tilde H_S,\cdot]+\Gamma_\downarrow\Dc[\sigma_-]+\Gamma_\uparrow\Dc[\sigma_+]+\tfrac12\Gamma_\varphi\Dc[\sz]$ with
$\Dc[L]X=LXL^\dagger-\tfrac12\{L^\dagger L,X\}$. Nonsecular Redfield for $\theta=\pi/2$ has
$\Gc_{xx}=0$, $\Gc_{yy}=-\Gamma_1$: starting from $\ket{+y}$ the Bloch vector initially shrinks along $y$
but not along $x$, which is fine; the violation appears through the $\Gc_{yx}=2\delta\omega$ term
combined with the inhomogeneity when $\delta\omega$ is comparable to $\Gamma_1$, i.e.\ when
$\gamma\gtrsim\omega_0$ so that $S(\omega)$ is broad.

\paragraph{Exercise~\ref{ex:5-4}.}
$K(0)=\Tr_B[\Lc_I\Qc\Lc_I(X\otimes\rhoB)]=\Tr_B[\Lc_I\Lc_I(X\otimes\rhoB)]$ since $\Pc\Lc_I\Pc=0$. Laplace:
invert \cref{eq:Dyson}. For $\theta=0$, $\hat\Pi_{01}(z)=\int_0^\infty\ee^{-(z+i\omega_0)t-\Gamma(t)}\dd t$, and
$\hat K_{01}(z)=z+i\omega_0-1/\hat\Pi_{01}$; expanding $\ee^{-\Gamma}=1-\Gamma+\Gamma^2/2-\dots$ with $\Gamma=O(g^2)$ produces
all powers of $g^2$ in $\hat K$: the irreducible crossed and nested diagrams do not cancel.

\paragraph{Exercise~\ref{ex:6-1}.}
$\Sigma(t)=\int_0^t\hat G(t,s)\Qc\tilde\Lc_I(s)\Pc\tilde G(s,t)\dd s$ is explicitly $O(\Lc_I)$ and vanishes at $t=0$.
Expanding $\hat G$, $\tilde G$ and $[1-\Sigma]^{-1}$ and using $\Pc\tilde\Lc_I\Pc=0$ gives
\cref{eq:TCL2} at second order and \cref{eq:TCL4} at fourth, where the subtracted terms
come from $\Pc\tilde\Lc_I\Sigma^2\Pc$ and from the expansion of $\tilde G(s,t)$ inside $\Sigma$. Odd
orders have an odd number of $B$'s.

\paragraph{Exercise~\ref{ex:6-2}.}
$\int_0^tS(\tau)\dd\tau=g^2\,\Re\big[(1-\ee^{-(\gamma/2-i\Omega)t})/(\gamma/2-i\Omega)\big]$; negative windows exist when
$\Omega>\gamma/2$, the first around $\Omega t\in(\pi,2\pi)$ for small $\gamma/\Omega$.

\paragraph{Exercise~\ref{ex:6-3}.}
With $\Pi^{(2)}(t)=\int_0^t\dd t_1\int_0^{t_1}\dd t_2\,\Pi_0(t-t_1)K^{(2)}(t_1-t_2)\Pi_0(t_2)$ and
$\Pi^{-1}=\Pi_0^{-1}-\Pi_0^{-1}\Pi^{(2)}\Pi_0^{-1}+\dots$, the fourth-order term of \cref{eq:TCLNZ} is
$\int K^{(4)}(t-s)\Pi_0(s)\Pi_0(t)^{-1}+\int K^{(2)}(t-s)[\Pi^{(2)}(s)\Pi_0(t)^{-1}-\Pi_0(s)\Pi_0(t)^{-1}\Pi^{(2)}(t)\Pi_0(t)^{-1}]$. The
bracket, written as a time-ordered integral, is $-\int\int\int$ of the products
$\avg{\Lc_I\Lc_I}\avg{\Lc_I\Lc_I}$ in the crossed and nested time orderings, which are exactly the
last two subtractions of \cref{eq:TCL4}; the reducible product is already absent from $K^{(4)}$.

\paragraph{Exercise~\ref{ex:6-4}.}
$\hat\rho_{ee}(z)=\widehat{|c_e|^2}(z)\rho_{ee}(0)$ and $\hat K_{ee}(z)=z-1/\widehat{|c_e|^2}(z)$; the Laplace
transform of $|c_e|^2$ is a rational function with poles at $-\gamma/2\pm i|d|/2$ and $-\gamma/2$,
so $\hat K$ is analytic to the right of them. The generator's rate $\gamma(t)$ has poles on the
real $t$ axis, so its Laplace transform does not exist as an ordinary function.

\paragraph{Exercise~\ref{ex:7-1}.}
Substitute $s^\pm=s^\cl\pm\tfrac12s^\qu$ and $C=S-\tfrac i2\chi$:
$(s^+_1-s^-_1)[Cs^+_2-C^*s^-_2]=s^\qu_1[S(s^+_2-s^-_2)-\tfrac i2\chi(s^+_2+s^-_2)]=s^\qu_1[Ss^\qu_2-i\chi s^\cl_2]$.
Induction: the sum over pairings of $2n$ commuting vertices is $\tfrac1{n!}(\text{one line})^n$.

\paragraph{Exercise~\ref{ex:7-2}.}
For a blip $s^\qu=\pm2$ on $[t_a,t_b]$, $-\int\!\!\int s^\qu Ss^\qu=-4\int_{t_a}^{t_b}\dd t_1\int_{t_a}^{t_1}\dd t_2\,S(t_1-t_2)=-Q_2(t_b-t_a)$
using $\ddot Q_2=4S$, $Q_2(0)=\dot Q_2(0)=0$. The blip--sojourn term with the preceding sojourn
$s^\cl=\pm1$ on $[t_0,t_a]$ gives $\pm i\cdot2\int_{t_a}^{t_b}\dd t_1\int_{t_0}^{t_a}\dd t_2\,\chi(t_1-t_2)$, which for
$t_0\to-\infty$ is $\mp iQ_1(t_b-t_a)$ with $\ddot Q_1=-2\chi$.

\paragraph{Exercise~\ref{ex:7-3}.}
$B_+(\tau)=\prod_kD(\alpha_k\ee^{i\omega_k\tau})$ with $\alpha_k=2g_k/\omega_k$. $D(\alpha_\tau)D(-\alpha)=\ee^{-i\alpha^2\sin\omega\tau}D(\alpha(\ee^{i\omega\tau}-1))$ and
$|\alpha(\ee^{i\omega\tau}-1)|^2=2\alpha^2(1-\cos\omega\tau)$, so $\avg{B_+(\tau)B_-(0)}=\ee^{-Q_2-iQ_1}$ and, with the
other ordering, $\avg{B_-(0)B_+(\tau)}=\ee^{-Q_2+iQ_1}$. Detailed balance: expand to first order
in $Q$, $\int_0^\infty\ee^{-i\varepsilon\tau}(-Q_2+iQ_1)\dd\tau$ picks $+2J(\varepsilon)[\coth(\beta\varepsilon/2)+1]/\varepsilon^2$, hence
$w_{\downarrow\leftarrow\uparrow}\propto1+n$, $w_{\uparrow\leftarrow\downarrow}\propto n$.

\paragraph{Exercise~\ref{ex:7-4}.}
(a) $\hat K^s(z)$ for the Ohmic kernel is
$\Delta^2\int_0^\infty\ee^{-z\tau}(1+\omega_c^2\tau^2)^{-\alpha}\cos(2\alpha\arctan\omega_c\tau)\,\dd\tau$; for $z\ll\omega_c$ it behaves as
$\hat K^s(z)\simeq\Delta^2\,\Gamma(1-2\alpha)\cos(\pi\alpha)\,z^{2\alpha-1}\omega_c^{-2\alpha}$, giving poles of $\hat P$ at
$z\sim\Delta_r\,\ee^{\pm i\pi/(2-2\alpha)}$ with $\Delta_r=\Delta(\Delta/\omega_c)^{\alpha/(1-\alpha)}\times O(1)$; for
$\alpha\ge\tfrac12$ the poles are on the negative real axis (incoherent
decay)~\cite[Sec.~V]{Leggett1987}. (b) See Example~7.2.

\paragraph{Exercise~\ref{ex:8-1}.}
Apply the identity to $i\mathsf LA$: $\dot A(t)=\ee^{i\mathsf Lt}i\mathsf LA=\ee^{i\mathsf Lt}[\mathsf P+\mathsf Q]i\mathsf LA$; the
$\mathsf P$ part is $i\Omega A(t)$, the $\mathsf Q$ part is expanded with the identity into $F(t)$ plus
$\int_0^t\ee^{i\mathsf L(t-s)}i\mathsf P\mathsf LF(s)\dd s$, and $i\mathsf P\mathsf LF(s)=A(A,A)^{-1}(A,i\mathsf LF(s))=-A(A,A)^{-1}(i\mathsf LA,F(s))=-AM(s)$
using self-adjointness of $\mathsf L$ and $(A,F)=0$.

\paragraph{Exercise~\ref{ex:8-2}.}
Project on the \emph{fluctuation} $A=\delta\sz=\sz-\avg{\sz}$, with
$(\delta\sz,\delta\sz)=1-\tanh^2(\beta\omega_0/2)$ (projecting on $\sz$ itself keeps the non-decaying
part $\avg{\sz}^2$ in $C(t)$ and gives a memory function that vanishes at $z=0$). Then
$i\mathsf L\delta\sz=i\sin\theta[\sx,\sz]B=2\sin\theta\,\sy B=F(0)$, $F(t)=2\sin\theta\,\sy(t)B(t)$ with
$\sy(t)=\sy\cos\omega_0t+\sx\sin\omega_0t$, and $\avg{\sy\sy(t)}=\cos\omega_0t+i\tanh(\beta\omega_0/2)\sin\omega_0t$,
$\avg{\sy(t)\sy}=$ its conjugate. Hence
\begin{equation*}
  (F,F(t))=2\sin^2\theta\big[\avg{\sy\sy(t)}C(-t)+\avg{\sy(t)\sy}C(t)\big]
  =4\sin^2\theta\Big[S(t)\cos\omega_0t-\tfrac12\tanh\tfrac{\beta\omega_0}{2}\,\chi(t)\sin\omega_0t\Big],
\end{equation*}
and $\hat M(0)=2\sin^2\theta[S(\omega_0)-\tanh(\beta\omega_0/2)\chi''(\omega_0)]/(1-\tanh^2)=2\sin^2\theta S(\omega_0)=\Gamma_1$
by the FDT. The response kernel enters with the system's polarization and is what makes the
normalization work out; there is no inhomogeneous term because the equilibrium value
sits in the reference state and the GLE describes the fluctuation around it.

\paragraph{Exercise~\ref{ex:8-3}.}
With $E_\mu=\sigma_\mu/\sqrt2$, $(E_\mu,\Lc^\dagger E_\nu)_B=\tfrac12\Tr[\sigma_\mu(i[H,\sigma_\nu])\rhoB]$, which is the transpose
of $\tfrac12\Tr[\sigma_\nu(-i[H,\sigma_\mu])\rhoB]$: Heisenberg generators are the transposes of
Schr\"odinger ones in a real orthonormal basis. The same holds for the kernel.

\paragraph{Exercise~\ref{ex:8-4}.}
In the RWA, $\mathsf Q i\mathsf L$ acting on $F\propto\sigma_-a^\dagger$-type operators closes on
$\sigma_-a$ (one excitation), and the second-level memory is the bath's own correlator
$\ee^{-(i\Omega+\gamma/2)t}$. The continued fraction $\hat C(z)=\tfrac12/[z+i\omega_0+g^2/(z+i\Omega+\gamma/2)]$ has the poles
of $c_e(t)$.

\paragraph{Exercise~\ref{ex:9-1}.}
Each $h^\cl$ derivative inserts an $S^\cl$ vertex as the \emph{last} operation before the trace
in at least one term; by the trace identity the full sum vanishes.

\paragraph{Exercise~\ref{ex:9-2}.}
$\sz=-i\eta_x\eta_y$, etc.; $\Theta\eta_x=-i\eta_x\eta_y\eta_z\eta_x=-i\eta_y\eta_z=\sx$. $H_I=S\otimes B$ becomes
$\Theta(\cos\theta\eta_z+\sin\theta\eta_x)B$; since $\Theta$ is conserved, in each sector $\Theta=\pm1$ the vertex is
linear in one Majorana and in $B$: the second-order self-energy of $\eta_i$ is a bath line
times a $\delta_{ij}$ Majorana loop-free diagram, i.e.\ $\Sigma_{\eta}\propto D$.

\paragraph{Exercise~\ref{ex:9-3}.}
Follow~\cite[Sec.~2]{Sieberer2016} with fermionic (Majorana) fields; the result is
$\Sigma^R=-i(\Gamma_\downarrow+\Gamma_\uparrow)/2$ and $\Sigma^K=-i(\Gamma_\downarrow-\Gamma_\uparrow)$ for the $\sigma_-$ propagator up to
normalization, so that $\Sigma^K/(\Sigma^R-\Sigma^A)=(\Gamma_\downarrow-\Gamma_\uparrow)/(\Gamma_\downarrow+\Gamma_\uparrow)=\tanh(\beta\omega_0/2)$ iff
detailed balance holds.

\paragraph{Exercise~\ref{ex:10-1}.}
Insert $C(\tau)=\int\frac{\dd\omega}{2\pi}C(\omega)\ee^{-i\omega\tau}$ and use $2\Re\int_0^t\ee^{i(\omega_0-\omega)\tau}\dd\tau=2\sin[(\omega-\omega_0)t]/(\omega-\omega_0)$.
A modulation $\omega_0\to\omega_0+\delta\cos\nu t$ replaces the sinc by a sum of sincs centred at $\omega_0+k\nu$
weighted by Bessel functions. If $C(\omega)$ grows above $\omega_0$ the broadened filter samples
more spectral weight and the decay accelerates (anti-Zeno).

\paragraph{Exercise~\ref{ex:10-2}.}
$\int_0^ty\ee^{i\omega s}\dd s=-\ee^{i\omega t/2}[2\cos(\omega t/2)-2]/(i\omega)$, so $|\cdot|^2=16\sin^4(\omega t/4)/\omega^2$. For the comb,
expand the square wave, $y=\tfrac4\pi\sum_{k\ \rm odd}k^{-1}\sin(k\pi s/\tau)$, and use
$|\int_0^t\ee^{i(\omega-k\pi/\tau)s}\dd s|^2\to2\pi t\,\delta(\omega-k\pi/\tau)$. The first tooth hits $\Omega$ when
$n\pi/t=\Omega$, i.e.\ $t=n\pi/\Omega$ ($=4\pi$ for $n=4$).

\paragraph{Exercise~\ref{ex:10-3}.}
Maximize $t\ee^{-2\Gamma_2t-2\kappa t^2}$: the quadratic term matters when $\kappa t^*\gtrsim\Gamma_2$, i.e.\
$t^*\lesssim\tau_B$; for GHZ states $t^*\sim1/(2N\Gamma_2)$, so the Zeno regime is reached for
$N\gtrsim1/(\Gamma_2\tau_B)$. For $T_1$-limited sensing (signal $\propto\ee^{-\Gamma_1t}$ with the phase on the
populations) the same algebra applies with $\Gamma_1$ in place of $2\Gamma_2$.

\paragraph{Exercise~\ref{ex:10-4}.}
Implement the pulse as $\rho\to\sx\rho\sx$ on the qubit factor at $t/2$ in the pseudomode
simulation. Expect (b) to be closer than (a), with a residual oscillating discrepancy of
relative size $\sim g^2\sin\theta\cos\theta/\omega_0$ at short times and a drift of the echo
``baseline'' from the $\chi$-driven inhomogeneity.
